\subsection{Three concurrent distinct leading lines}
\label{sec:concurrentregular}

Let the common vertex be the $T$ direction, so the leading form
is a squarefree binary cubic $F(u,v)$. Its binary gradient has
only the zero at the origin. If $F_2$ at the vertex is nonzero,
a sequence with $\nabla P=o(\|x\|)$ must approach that vertex;
then $u,v=o(T)$ and the $T$ equation has a nonzero linear
leading term. The Gaussian gradient criterion excludes purity.
An isotropic zero-$F_2$ vertex already gives a frame.

In the remaining alternative retain the full Gaussian with
$m=q_{TT}\ne0$ and write
\[
 P=T H(u,v)+G(u,v),\qquad\deg H\le1,\quad G_3=F,
 \qquad q=mT^2/2+T J(u,v)+q_*(u,v),
\]
where $J$ is affine linear and $q_*$ includes all lower terms.
For each fixed $t$, the phase gradient is linearly coercive:
its only possible limiting singular direction is the vertex,
where its $T$ derivative has leading term $-mT$. The estimate
is locally uniform on compact parameter sets. Completing $T$
exactly gives
\[
 tP-q=-\frac m2\left(T-\frac{tH-J}{m}\right)^2
                    +tG-q_*+\frac{(tH-J)^2}{2m}.
\]
This change has volume Jacobian one. The binary phase has
squarefree cubic leading form $tF$, whose two leading quadratic
derivatives are a regular sequence of length four. Weighted
leading reduction, or the same Koszul argument used above,
gives full twisted rank four. The Gaussian factor has rank one.
Condition~\eqref{eq:nonisoleading} holds at the sole leading
singular direction. Thus all escapes are excluded from every
pure linear-action block by Proposition~\ref{prop:boundedregularrank}.

If $H$ is a nonzero constant, $P$ has no critical point.
The full regular rank is zero and a constant original marked
period is impossible. This remains true after a constant shift
of $P$.

If $H=0$, integrate the $T$ Gaussian formally. A fixed polynomial
marker becomes a fixed polynomial in $u,v$, using conditional
mean $-J/m$ and covariance $1/m$. The effective Gaussian is
\[
 q_{\rm eff}=q_*-J^2/(2m).
\]
Its Hessian is nondegenerate by the Schur determinant identity
$\det\operatorname{Hess}q=m\det\operatorname{Hess}q_{\rm eff}$.
The squarefree leading binary cubic gives
$\|\nabla G\|\ge c\|(u,v)\|^2$ outside a compact set.
The all-degree Gaussian gradient criterion therefore excludes
every nonzero finite exponential-polynomial original period.
The conditional normalization is independent of $t$.

Finally let $H$ be nonconstant and use $u=H$ as an affine
coordinate. Then $P=uT+G(u,v)$. If $R(v)=G(0,v)$ is
nonconstant, its critical points give all the bounded groups.
The exact smooth-constraint shear gives their local types
$\alpha+s z+y^{\mu+1}$. This includes the empty critical set
when $R$ is linear. Their full power characters have positive
original-marker orders, so the regular-rank and seed lemmas
exclude a constant period without constructing the escaping
blocks.

If $R=k$ is constant, write $G=k+uB(u,v)$, where $B$ has
degree at most two. The critical curve is
\[
 u=0,\qquad T=-B(0,v).
\]
The quadratic coefficient $c$ of $B(0,v)$ is nonzero.
Otherwise the squarefree cubic $F$ would be divisible by $u^2$.
Put $s=T+B(u,v)$, so exactly $P=k+us$ and the coordinate
Jacobian is one. Scale $u=\rho U$, $s=\rho V$,
$\rho=t^{-1/2}$. The fixed phase is
\[
 UV-q(0,v,-B(0,v)),
\]
with quartic coefficient $-mc^2/2\ne0$ in $v$. Its rank and
critical Milnor number are three. The exact critical equations
$u=\epsilon q_s$, $s=\epsilon q_u$ show that every bounded
branch has bounded scaled coordinates. This is precisely the
complete quartic parabola construction above, now with
Jacobian $\rho^2$. It supplies the entire regular block when
$k=0$, with positive original-marker orders; when $k\ne0$
that block has nonzero exponential part. The polynomial
representatives of $v$ powers are original markers, so its
leading matrix is invertible. The seed lemma applies.

These alternatives exhaust the affine polynomial $H$ and the
restriction $R$. They also exclude nonzero finite
exponential-polynomial markers after differential isolation and
shifting the constant term. We obtain the following conclusion.

\begin{proposition}\label{prop:concurrentcomplete}
Every pure cubic whose leading form is three distinct concurrent
lines has an isotropic vertex with $F_2=0$. Hence it has a
Gaussian affine-linear coordinate and satisfies $m>3\deg Q$.
\end{proposition}
\begin{proof}
The gradient gate and the nonisotropic cases just proved leave
only that vertex alternative. Apply Lemma~\ref{lem:frame} and
the complete affine-linear theorem.
\end{proof}
